\chapter{Security Model and the Version 8 Hardening}\label{ch:security}

A full review of protocol version~7 in October 2026 found 31 issues of security, correctness and performance,
numbered \#1--\#31 in the repository's roadmap~\cite{repo}. This chapter states the threat model, explains the
most important findings in detail (what each allowed, how it was fixed, and why that fix was chosen over the
alternatives), and then summarises every fix in one table. Each fix was made as its own commit with its own
tests. The socket-style SDK and the HTTP mapping were deliberately excluded and are discussed in
Chapter~\ref{ch:conclusion}.

\section{Threat model}

We consider three attackers:
\begin{enumerate}
  \item An \textbf{off-path attacker}, who can send datagrams with arbitrary, spoofed source addresses but cannot
    see traffic between the endpoints.
  \item An \textbf{on-path passive attacker}, who can observe all traffic.
  \item An \textbf{on-path active attacker}, who can additionally drop, delay, reorder, modify, replay and inject
    datagrams.
\end{enumerate}

In \emph{plaintext mode}, Fuse makes no confidentiality or integrity claim against an on-path attacker. It must
still resist off-path attackers: they must not be able to hijack a session, inject acknowledgements that cause
data loss, use an endpoint as a traffic amplifier, or crash or exhaust it with malformed input.

Under a \emph{pre-shared key}, Fuse additionally aims for confidentiality and integrity of all data, and
integrity of all control signals, against an active on-path attacker. Denial of service by an on-path attacker
is out of scope, since such an attacker can always drop everything. So is forward secrecy, which the PSK path
explicitly does not provide.

\section{AES-GCM nonce reuse on retransmission (\#1)}\label{sec:nonce}

\paragraph{Problem.} The AAD of an encrypted DATA block included its flags byte, and a retransmission sets the
retransmission bit in that byte. A retransmitted block was therefore sealed under the \emph{same} key and the
\emph{same} nonce $(\mathit{lane},\mathit{seq})$ as the original, but with a \emph{different} AAD.

\paragraph{Consequence.} GCM's authentication tag is a polynomial in the secret hash key $H$ evaluated over the
AAD and ciphertext. Two tags computed under one (key, nonce) pair over different inputs give an attacker two
equations in $H$, from which $H$ can be recovered. With $H$ known, the attacker can forge valid tags for
arbitrary blocks: this is the ``forbidden attack'' on GCM~\cite{joux,nist-gcm}. Every retransmission, which
happens routinely on any lossy path, handed a passive observer such a pair.

\paragraph{Fix.} The retransmission bit is masked out of the AAD (\code{build_aad}, \code{transfer.cpp:105}). A
resend is now byte-for-byte the original ciphertext and tag, which reveals nothing new. The last-block bit
remains authenticated, so it still cannot be flipped to end a lane early.

\paragraph{Why this fix.} An alternative was to give every transmission a fresh nonce, as QUIC does by never
reusing packet numbers. That requires a separate packet number on the wire and a mapping from packet numbers to
blocks on both ends, a much larger change. The retransmission bit is only a hint used to measure RTT, so it does
not need integrity protection.

\section{Unauthenticated control plane under a pre-shared key (\#2)}

\paragraph{Problem.} Only DATA payloads were protected. ACK, NACK, StreamStart, ResumeRanges, StreamClose and
FileDigest were plaintext and unauthenticated even when encryption was enabled, so an active attacker could
manipulate the control plane freely:
\begin{itemize}
  \item a forged ACK claims delivery of data the receiver never got, so the sender releases it and the transfer
    ends incomplete;
  \item a forged NACK halves the window;
  \item a forged close ends a lane;
  \item a forged ResumeRanges message makes the sender skip data.
\end{itemize}

\paragraph{Fix.} Every control datagram on the bulk-transfer path now carries a 24-byte trailer: a 64-bit
counter and a GMAC tag over the whole datagram. The tag is computed under a \emph{control key} derived
separately from the data keys, with a nonce space disjoint from that of the data (Section~\ref{sec:wirecrypto}).
Unsigned control messages are discarded before they are parsed.

\paragraph{Why this fix.} GMAC reuses the AES-GCM context the lane already has, costs one AES pass over a few
dozen bytes, and adds no round trips. Running the control plane through DTLS would have reintroduced the
record-layer state that the design avoids (Section~\ref{sec:crypto}). The counter need only be unique per
sender, not checked for replay. Replaying an authentic control message is harmless: acknowledgement bases only
move forward, and an abort can only be replayed if it was really sent.

\paragraph{Exceptions.} Three messages remain unauthenticated by necessity. In each case forgery can only cause
a denial of service that an on-path attacker could cause anyway.
\begin{itemize}
  \item \emph{Retry} is sent before any key is confirmed. A forged one can only make the token wrong, and the
    receiver then answers again.
  \item The \emph{version-mismatch reply} must be understood by any version.
  \item \emph{StreamClose(auth failed)}: the receiver shares no key with a mis-keyed sender.
\end{itemize}

\section{Key confirmation (\#6)}

\paragraph{Problem.} With mismatched pre-shared keys, a transfer only failed after twenty data blocks had failed
authentication, or after a full timeout.

\paragraph{Fix.} The StreamStart is now signed with a key derived from its own salt, so the receiver confirms
the key before allocating anything. On failure it answers StreamClose(reason~2), and the sender reports
\code{AuthFailed} at the handshake. The same check catches a plaintext receiver facing an encrypting sender, and
the reverse.

\section{Address validation and amplification (\#5, \#18)}

\paragraph{Problem.} Two weaknesses combined. First, the first StreamStart to reach a receiver lane claimed it,
whatever its source address, so an off-path attacker could race a legitimate sender. Second, the receiver
answered StreamStarts with acknowledgements and ResumeRanges messages larger than the request, and kept sending
periodic acknowledgements to a peer that never answered. Both made it a reflection amplifier: spoof a victim's
address, and the receiver sends traffic to the victim.

\paragraph{Fix.} The fix has two parts.
\begin{itemize}
  \item A StreamStart without a valid \emph{address token} receives only a 20-byte Retry, which is smaller than
    the 80-byte request. The token is SipHash~\cite{siphash} of the source address and lane nonce under a
    secret held by the receiver lane. Nothing is allocated or claimed until a StreamStart echoes the token, which
    only a sender that receives at that address can do.
  \item Timer-driven acknowledgements stop after three sent to a peer that has fallen silent.
\end{itemize}

\paragraph{Why this fix.} It is the mechanism QUIC uses for the same problem~\cite{rfc9000}, and it is
stateless on the receiver: the token can be recomputed from the datagram, so validating an address needs no
memory per sender. A cookie exchange that stored state per sender would itself be a resource-exhaustion
target.

\section{Path validation (\#3)}

\paragraph{Problem.} To survive NAT rebinding, both ends re-anchored their peer address to the source of any
accepted datagram. A single spoofed datagram could therefore redirect a lane's traffic: a hijack, or a
reflection at a third party.

\paragraph{Fix.} A new source address is adopted only after it echoes a random 64-bit challenge in a
PathResponse (Section~\ref{sec:path}). Data from the new source is accepted in the meantime, so a genuine
rebinding costs no throughput. The design follows QUIC's path validation.

\section{Predictable nonces (\#20)}

\paragraph{Problem.} Lane nonces and session nonces came from a seeded \code{mt19937_64}. Its future outputs can
be predicted from a few observed ones. Since nonces were the only session check on several control messages, an
attacker who observed a few could forge messages for later sessions.

\paragraph{Fix.} \code{secure_random_u64()} draws from the cryptographic library's random generator, or from the
kernel's \code{getrandom}/\code{getentropy} in builds without one, and aborts if no secure source exists.

\section{Key usage limit (\#21)}

\paragraph{Problem.} One AES-GCM key protected every block of a lane, without limit. RFC~9001 applies a
confidentiality limit of about $2^{23}$ packets per key~\cite{rfc9001}, which a multi-gigabyte transfer
approaches.

\paragraph{Fix.} Data keys are derived per lane and per epoch of $2^{22}$ blocks, identically on both ends, with
no messages on the wire (Section~\ref{sec:crypto}).

\paragraph{Why this fix.} An explicit rekey exchange would need a new message, acknowledgement and transition
logic, and its own failure modes. A deterministic schedule needs none of these, because both ends already agree
on the sequence number of every block.

\section{Weak pre-shared keys (\#19, \#28)}

The pre-shared key is used directly as HKDF input, without stretching, so a human-chosen password could be
guessed offline from one captured handshake at full speed. The bulk transfer now refuses keys shorter than 16
bytes, and the documentation requires random keys. Accepting passwords safely would require a
password-authenticated key exchange (PAKE), which remains open. Separately, keys longer than 256 bytes used to
fail key derivation because of a fixed-size buffer; the key is now copied into a buffer of its own length.

\section{Information leaks (\#31)}

Two leaks were fixed:
\begin{itemize}
  \item DATA carried the raw monotonic clock, which reveals the host's uptime. It now carries the clock plus a
    random per-stream offset (Section~\ref{sec:wirecrypto}).
  \item The DTLS wrapper kept a debug copy of all wire traffic in the production path; this is now opt-in.
\end{itemize}

\section{Transport correctness and performance fixes}

The remaining findings concerned the transport's behaviour rather than its security. In summary:
\begin{itemize}
  \item \textbf{\#10 SACK.} The bitmask in every acknowledgement was ignored. \code{apply_ack} now releases
    blocks acknowledged selectively, and the timeout resends known holes (Sections~\ref{sec:registry}
    and~\ref{sec:recovery}).
  \item \textbf{\#23 Timer.} A fixed multiple of SRTT fired spuriously under jitter and never backed off. It is
    replaced by $\max(\mathit{SRTT}+4\,\mathit{RTTVAR},2\,\mathit{SRTT})$ with exponential backoff and one
    halving per run of timeouts.
  \item \textbf{\#8 Window and slow start.} The window counted blocks, so it inflated up to 13 times as blocks
    grew. It now counts units of the configured block size, and the transfers use slow start.
  \item \textbf{\#22 Pacing.} Windows left in line-rate bursts. Sending is now paced over the round trip, with
    smaller bursts on real paths.
  \item \textbf{\#9 Flow control.} There was none. Acknowledgements now carry a receive window, with window
    updates (Section~\ref{sec:flow}).
  \item \textbf{\#7 MTU.} Blocks grew to 16~KiB regardless of the path, causing IP fragmentation. Growth is now
    capped at the route MTU, with a ceiling after loss and later re-probing (Section~\ref{sec:blocksize}).
  \item \textbf{\#16 and \#17 Algorithmic bounds.} Receiving was $O(W)$ and a reassembly set was unbounded. Both
    now use rings bounded by the window.
  \item \textbf{\#11 Versions and \#24 extensions.} Version mismatches were silent and every new field needed a
    version bump. Version~8 adds the explicit mismatch reply and TLV extension areas.
  \item \textbf{\#15 File identity.} A file rewritten in place could be resumed with stale data. The identity now
    includes the inode change time, inode number and device.
  \item \textbf{\#25 Ports.} \code{base_port + lane} wrapped past 65,535; this is now a configuration error.
  \item \textbf{\#30 Validation.} SETUP was not validated, the multiplexed transfer indexed streams with
    unchecked identifiers, NACKs could be duplicated or dropped, and \code{SendQueue(0)} divided by zero. All are
    fixed.
  \item \textbf{IPv6.} Senders always opened IPv4 sockets, so IPv6 hosts were unreachable. Sockets are now opened
    in the destination's family.
\end{itemize}

\section{Summary of all fixes}

\begingroup\small
\begin{xltabular}{\linewidth}{c>{\raggedright\arraybackslash}p{4.6cm}>{\raggedright\arraybackslash}X}
\caption{Every issue fixed in the protocol core for version~8.}\label{tab:fixes}\\
\toprule
\textbf{\#} & \textbf{Issue} & \textbf{Fix} \\
\midrule\endfirsthead
\toprule
\textbf{\#} & \textbf{Issue} & \textbf{Fix} \\
\midrule\endhead
\bottomrule\endfoot
1 & GCM nonce reuse on every retransmission & Retransmission bit masked out of the AAD; a resend is the identical ciphertext and tag \\
2 & Control plane unauthenticated under a PSK & Counter plus GMAC trailer under a separate control key; unsigned messages ignored \\
3 & Address migration without validation & PathChallenge/PathResponse before the peer address changes \\
5 & No session binding without encryption; replayable SETUP & Retry tokens bind a lane to a validated address; SETUP responder nonce echoed in FINACK \\
6 & No key confirmation & Signed StreamStart; mismatch reported as \code{AuthFailed} at the handshake \\
7 & Block growth unaware of the MTU & Growth capped at the route MTU; ceiling after loss; DPLPMTUD-style re-probing \\
8 & Window counted in blocks; no slow start & Window in configured-block units, capped by the largest permitted block; opt-in slow start \\
9 & No flow control & \code{rwnd} in every ACK; the sender stays within it; window updates \\
10 & SACK bitmask discarded & \code{apply_ack} releases selectively acknowledged blocks; RTO resends known holes \\
11 & Silent version mismatch & Fixed-layout reply; \code{VersionMismatch} status \\
15 & File identity from name, size and mtime only & Also hashes ctime, inode and device \\
16 & $O(W)$ receive path & Ring-indexed receiver window \\
17 & Unbounded set in reassembly & Window-bounded ring; capped trace \\
18 & Amplification from the receiver & Retry before anything else; idle ACKs stop after three \\
19 & PSK brute-forceable if a password & Keys under 16 bytes refused; random keys required (PAKE open) \\
20 & Predictable nonces & Cryptographic RNG or \code{getrandom}; abort if neither \\
21 & No AES-GCM usage limit & Lane keys rotate every $2^{22}$ blocks, derived on both ends \\
22 & No pacing & Pacing over the RTT; smaller bursts on real-RTT paths \\
23 & Aggressive RTO without backoff & $\max(\mathit{SRTT}+4\mathit{RTTVAR},2\mathit{SRTT})$, exponential backoff, one halving per run \\
24 & No extension mechanism & TLV areas in StreamStart and SETUP with critical-bit semantics \\
25 & Port number wrap-around & \code{ConfigError} \\
28 & PSK over 256 bytes failed (core part) & Buffer sized to the key \\
30 & Unvalidated SETUP; mux indexing; dropped NACKs; \code{SendQueue(0)} & Responder validation; identifiers $0..n-1$; NACK de-duplication with immediate resend; capacity clamp \\
31 & Debug wire copy in DTLS; uptime in \code{send_time_ns} & Opt-in recording; random per-stream clock offset \\
--- & IPv4-only sender sockets & Socket opened in the destination's family \\
\end{xltabular}
\endgroup

\section{Security properties after version 8}

Under a pre-shared key, the bulk transfer now provides the following properties:
\begin{itemize}
  \item confidentiality and integrity of data, with no nonce reuse;
  \item integrity of every control decision the protocol acts upon;
  \item key confirmation at the handshake;
  \item bounded key usage;
  \item resistance to off-path session hijacking, address redirection and reflection amplification.
\end{itemize}

In plaintext mode, address validation, path validation, unpredictable nonces and strict input validation still
defeat the off-path attacks, but an on-path attacker can read and modify everything, which is inherent in
plaintext.

The pre-shared-key path still does not provide forward secrecy, and the socket-style SDK and the HTTP mapping do
not yet have these protections (Chapter~\ref{ch:conclusion}).
